\textbf{Data and tasks.} Breast cancer Wisconsin (diagnostic) from scikit-learn: 569 patients, 30 features, 212 malignant (positive) and 357 benign. Task \texttt{bc\_natural} uses the data as is. For ratios $k\in\{5,20,50\}$ (\texttt{bc\_1to5}, \texttt{bc\_1to20}, \texttt{bc\_1to50}) we make the \emph{training} set imbalanced: after a stratified 70/30 split (seeded), training positives are subsampled to $\mathrm{round}(n_0/k)$, which gives 50, 12 and 5 positives in every split (the natural split has 148). Subsampling the whole data to 1:50 would leave about 7 positives in total, so every held-out example is kept (64 malignant, 107 benign) and test positives get weight $w=(n_0^{\text{test}}/k)/n_1^{\text{test}}$ so that the \emph{weighted} test prevalence equals the training prevalence $1:k$; all metrics are weighted. This tests calibration at the deployment prevalence but reduces the effective number of positives at 1:50 (Limitations).

\textbf{Systems and protocol.} Five strategies $\times$ two learners $\times$ four tasks $\times$ 10 seeds (0--9; each seed changes the split, the positive subsample and the resampling randomness), 400 runs of \texttt{method/run.py} with identical code, data pipeline and metric definitions. Hyperparameters are fixed (Section~\ref{sec:method}); no system is tuned, so the tuning budget is equal (zero). Correction target ratio is 1:1 unless stated. Ablations: the prior-shift offset on reweighting, ROS and SMOTE for both learners and all tasks (240 runs), and a sweep of the SMOTE target ratio $r\in\{0.1,0.25,0.5,0.75\}$ plus the main runs at $r=1$, both learners, at 1:20 (80 runs). Reference arm (``method'' in the run registry): LR with no correction.

\textbf{Statistics.} Differences to the same learner without correction are paired by seed, with a two-sided paired $t$-test; $p$-values are unadjusted. They are computed by \texttt{analysis/make\_tables.py} (SciPy), written to \texttt{results/tables/*.csv}, and agree to machine precision with \texttt{rh compare} (\texttt{results/tables/compare/}). Isolated $p<0.05$ cells are suggestive only. Calibration slope is heavy-tailed (the registry holds values from 0.19 to 18.6), so slopes and $|b-1|$ are reported as medians with interquartile ranges; all other metrics as means.

\textbf{Hypotheses (registered in \texttt{proposal.md} before the main runs).} H1: reweighting, ROS and SMOTE give no AUROC/AUPRC gain of at least 0.01 over no correction (refuted for a cell if the mean paired gain is $\ge0.01$ with $p<0.05$). H2: they raise Brier score and $|a|$ at 1:20 and 1:50. H3: threshold moving leaves probability metrics identical and trades specificity for sensitivity. H4: with the prior-shift offset, Brier score of the corrected systems is within 0.005 of no correction.

\textbf{Hardware and runtime.} CPU only (2 threads), each run takes about 2\,s including process start-up; the 400 main runs took about 7 minutes of wall-clock with two processes. One early sanity run failed (missing output directory) and was rerun; it stays in the registry.
